% Generated draft: requires mathematical and source review.
\documentclass[11pt]{article}
\usepackage[margin=25mm]{geometry}
\usepackage{fontspec}
\setmainfont{Latin Modern Roman}
\usepackage{amsmath,amssymb,mathtools,unicode-math}
\setmathfont{Latin Modern Math}
\usepackage{enumitem,needspace,xurl,hyperref}
\setlength{\parindent}{0pt}
\setlength{\parskip}{5pt}
\setlength{\emergencystretch}{3em}
\newcommand{\N}{\mathbb{N}}
\newcommand{\Z}{\mathbb{Z}}
\newcommand{\Q}{\mathbb{Q}}
\newcommand{\R}{\mathbb{R}}
\newcommand{\C}{\mathbb{C}}
\newcommand{\F}{\mathbb{F}}
\newcommand{\A}{\mathbb{A}}
\newcommand{\PP}{\mathbb{P}}
\newcommand{\E}{\mathbb{E}}
\newcommand{\eps}{\varepsilon}
\newcommand{\dd}{\,\mathrm{d}}
\begin{document}
\section*{UnsolvedMath ID: 30003722}
\textbf{Source problem code:} catalogue entry
\subsection{Definitions and mathematical statement}
Let \(G=(V,E)\) be an infinite graph. The supplied statement refers to first-passage percolation in a hostile environment, with competing types \(F_1\) and \(F_\lambda\), parameter \(\lambda\), and initial red-particle density \(\mu\). It does not specify the stochastic dynamics or probability law, so these types and parameters cannot be fully defined from the supplied record alone.

Conditional on a specification of that process, let \(V_1\) and \(V_\lambda\) denote the sets of vertices ultimately occupied by the respective types. A type cluster means a connected component of the induced subgraph \(G[V_i]\), for \(i\in\{1,\lambda\}\). Define the events
\[
\mathcal I_i=\{G[V_i]\text{ has an infinite connected component}\},
\qquad
\mathcal C=\mathcal I_1\cap\mathcal I_\lambda.
\]
This expresses coexistence as the simultaneous existence of an infinite cluster of each type in the same realization; the interpretation of the types as ultimate occupied vertex sets requires confirmation from the source.

The first question is to determine the graph and parameter conditions under which \(\mathcal I_\lambda\) occurs. The original formulation restricts attention to hyperbolic graphs, whereas the supplied clean statement asks this on bounded-degree graphs. These scopes are retained as an unresolved distinction.

The second question concerns \(G=\Z^d\): for every \(\lambda\in(0,1)\), does there exist an interval \(I_\lambda\) of admissible initial red-particle densities such that coexistence occurs for every \(\mu\in I_\lambda\)? If coexistence is intended to mean positive probability, the proposed assertion would be
\[
\forall\lambda\in(0,1)\quad
\exists I_\lambda\quad
\forall\mu\in I_\lambda:\qquad
\mathbb P_{\Z^d,\lambda,\mu}(\mathcal C)>0.
\]
Here the probability law is a placeholder for the missing process specification. The record does not determine whether positive probability or probability one is intended, which dimensions \(d\) are included, or whether the interval must have positive length and exclude endpoint densities. Consequently, this display records a conditional interpretation of the conjecture, not a fully specified assertion.
\subsection{Short English statement}
Determine when the hostile-environment process produces an infinite \(F_\lambda\) cluster, and whether, on \(\Z^d\), every \(\lambda\in(0,1)\) admits an interval of red-particle densities supporting infinite clusters of both types simultaneously. The process law, dimension range, interval requirements, probability quantifier, and graph scope remain insufficiently specified.
\subsection{Sources}
\begin{itemize}
\item Local UnsolvedMath registry, \texttt{data/unsolvedmath/problems.json}, permanent ID \(30003722\), source code \texttt{catalogue entry}; supplied title: ``Competing First-Passage Percolation in Hostile Environments.'' Collection: Oberwolfach Reports Open Problems. No individual proposer is supplied; the proposed year is recorded as 2018.
\item Supplied report citation: \emph{Strongly Correlated Random Interacting Processes} (2019), Oberwolfach Reports, DOI 10.4171/owr/2018/4, pp.\ 187--253. \url{https://doi.org/10.4171/owr/2018/4}. The literature-source label identifies this as ``Oberwolfach Report 4/2018''; the supplied bibliographic year and report-year designation require reconciliation.
\item Supplied literature source: Candellero--Stauffer, ``competing FPP.'' \url{https://arxiv.org/abs/1810.04593}. No fuller bibliographic description is supplied.
\item Literature assessment dated exactly \texttt{2026-08-21}: the supplied curation labels the problem partially solved and reports proofs of coexistence and an infinite \(F_\lambda\) cluster in substantial parameter regimes on vertex-transitive non-amenable hyperbolic graphs, while describing the broader conjectured parameter range and Euclidean cases as unresolved. This is supplied dated curation, not an independently verified assessment. The background expressly requests independent verification of time-sensitive status and cited claims.
\item The registry labels the statement \texttt{corrected\_verified} and reports a manual semantic, typographic, and LaTeX review asserting self-containment and preservation of the target. These are supplied review claims; they do not resolve the missing mathematical specifications identified here. No cited source was independently consulted for this draft.
\end{itemize}
\end{document}
